\section{Conclusion}
\label{sec:conclusion}

We asked what a retrieval-augmented time-series forecaster should retrieve, and
answered that it should retrieve its own errors---and that the futures its
predecessors retrieve are the same correction with one term forced on. \method{}
realizes that answer as a post-hoc layer over a chronological memory of a
forecaster's validation residuals, with shrinkage, a dial over the retrieval object,
and a gate that can leave the forecast untouched. With all 585 base forecasts fixed
it reaches 541 strict wins against 490 for pure error retrieval and 478 for the
analogue control, and takes the best average in every task family.

\paragraph{Scope and limitations.}
Error memory pays off when a forecaster's mistakes recur and borrowing a
neighbour's trajectory pays off when they do not; Figure~\ref{fig:object-strength}
locates that crossover by base accuracy and Section~\ref{sec:exp-analysis} by
horizon. Leaving the dial free costs 11 strict wins against pinning it shut, and a
horizon-dependent prior on it recovers those wins while giving back most of the
long-term magnitude, so which of the two to deploy is a choice we report rather
than settle. The end-to-end and additional-system results use the error-only
configuration, so what the dial adds is measured only with base forecasts fixed. The
primary matrix uses one seed and the prequential study one backbone; descriptor,
pooling, and shrinkage ablations remain open.\label{sec:main-text-end}
